% 霍普夫代数（综述）
% license CCBYSA3
% type Wiki

本文根据 CC-BY-SA 协议转载翻译自维基百科\href{https://en.wikipedia.org/wiki/Hopf_algebra}{相关文章}

在数学中，霍普夫代数（Hopf algebra，以海因茨·霍普夫 Heinz Hopf 命名）是一种结构，它既是一个（含幺结合的）代数，又是一个（含余元的余结合）余代数；这些结构的相容性使它成为一个双代数，此外它还配备了一个满足特定性质的反同态。霍普夫代数的表示论尤其优美，因为兼容的余乘法、余元以及反元的存在，使得能够构造表示的张量积、平凡表示以及对偶表示。

霍普夫代数自然地出现在代数拓扑中（其起源之处，并与 H-空间概念相关）、群概形理论、群论（通过群环的概念），以及许多其他地方，因此可能是最常见的一类双代数。霍普夫代数也作为独立的研究对象受到关注，一方面有许多关于具体类实例的研究，另一方面也存在分类问题。它们拥有广泛的应用，从凝聚态物理与量子场论，到弦论以及大型强子对撞机（LHC）现象学。
\subsection{形式定义}
形式上，霍普夫代数是域$K$上的一个（结合且余结合的）双代数$H$，并配备一个$K$-线性映射$S: H \to H$（称为反元，antipode），使得下列图表交换：
\begin{figure}[ht]
\centering
\includegraphics[width=8cm]{./figures/080fad97a4c9a060.png}
\caption{} \label{fig_HPFds_1}
\end{figure}
这里，$\Delta$是双代数的余乘法，$\nabla$是其乘法，$\eta$ 是其单位元映射，$\varepsilon$是其余元映射。在无求和的Sweedler 记号下，这一性质也可以表述为：
$$
S(c_{(1)})c_{(2)} \;=\; c_{(1)}S(c_{(2)}) \;=\; \varepsilon(c)\,1 
\qquad \text{对于所有 } c \in H.~
$$
与代数的情况类似，在上述定义中，可以将底层的域$K$替换为交换环$R$。\(^\text{[4]}\)

霍普夫代数的定义是自对偶的（这在上面图表的对称性中有所体现），因此如果可以定义 $H$的对偶（当$H$是有限维时总是可能的），那么该对偶也自动是一个霍普夫代数。\(^\text{[5]}\)
\subsubsection{结构常数}
在固定了底层向量空间的一组基$\{ e_k \}$之后，可以通过乘法的结构常数来定义代数：
$$
e_i \nabla e_j = \sum_k \mu_{ij}^k e_k~
$$
对于余乘法：
$$
\Delta e_i = \sum_{j,k} \nu_i^{\,jk} \, e_j \otimes e_k~
$$
以及反元：
$$
S e_i = \sum_j \tau_i^{\,j} e_j~
$$
结合律要求：
$$
\mu_{ij}^k \, \mu_{kn}^m \;=\; \mu_{jn}^k \, \mu_{ik}^m~
$$
而余结合律要求：
$$
\nu_k^{\,ij} \, \nu_i^{\,mn} \;=\; \nu_k^{\,mi} \, \nu_i^{\,nj}~
$$
连接公理要求：
$$
\nu_k^{\,ij} \, \tau_j^{\,m} \, \mu_{im}^n 
\;=\; 
\nu_k^{\,jm} \, \tau_j^{\,i} \, \mu_{im}^n~
$$
\subsubsection{反元的性质}
反元$S$有时被要求具有一个$K$-线性逆映射，这在有限维情形下会自动成立【需要澄清】，或者当$H$是交换的、余交换的（或更一般地是拟三角的 quasitriangular）时也成立。

一般而言，$S$是一个反同态\(^\text{[6]}\)，因此$S^2$是同态；若$S$可逆（通常要求如此），那么 $S^2$ 就是一个自同构。

若$S^2 = \mathrm{id}_H$，则称该霍普夫代数为对合的，其底层带对合的代数是一个 $*$-代数。若$H$是有限维、在零特征域上半单、交换或余交换的，那么它是对合的。

如果一个双代数$B$承认一个反元$S$，那么这个$S$是唯一的（“一个双代数至多承认一个霍普夫代数结构”）\(^\text{[7]}\)。因此，反元并不是可以随意选择的额外结构：作为霍普夫代数，是双代数的一个性质。

反元可以类比为群上的逆映射，即把元素$g$映射为$g^{-1}$\(^\text{[8]}\)。
\subsubsection{霍普夫子代数}
若$H$是一个霍普夫代数，$A$是其子代数，则当$A$同时也是$H$的子余代数，并且反元$S$将$A$映射到$A$中时，称$A$为一个霍普夫子代数。换言之，当把$H$的乘法、余乘法、余元和反元限制到$A$上时，$A$本身就是一个霍普夫代数（此外还要求$H$的单位元 $1$属于$A$）。Warren Nichols 与 Bettina Zoeller 于 1989 年证明的 **Nichols–Zoeller 自由性定理表明：如果$H$是有限维的，则自然的$A$-模 $H$ 是有限秩的自由模；这推广了子群的拉格朗日定理\(^\text{[9]}\)。由此及积分理论推出，半单有限维霍普夫代数的霍普夫子代数自动是半单的。

在霍普夫代数$H$中，若一个霍普夫子代数$A$满足右稳定性条件，即$\operatorname{ad}_r(h)(A) \subseteq A \quad \text{对所有 } h \in H$，则称$A$在$H$中是右正规的。其中右伴随映射$\operatorname{ad}_r$定义为$\operatorname{ad}_r(h)(a) = S(h_{(1)})\, a \, h_{(2)}, \quad a \in A, \; h \in H$.类似地，若$A$在左伴随映射$\operatorname{ad}_l(h)(a) = h_{(1)}\, a \, S(h_{(2)})$下稳定，则称$A$在$H$中是左正规的。当反元$S$是双射时，左右正规条件等价，此时称$A$为一个正规霍普夫子代数。

在$H$中的正规霍普夫子代数$A$满足以下条件（作为$H$的子集相等）：$H A^{+} = A^{+} H$,其中$A^{+}$表示$A$上余元的核。此正规性条件意味着$H A^{+}$是$H$的一个霍普夫理想（即：它既是余元核中的代数理想，也是余代数的余理想，并且在反元下稳定）。因此，可以得到一个商霍普夫代数$H / H A^{+}$,以及一个满射$H \to H / A^{+} H$,其理论与群论中正规子群和商群的理论类似\(^\text{[10]}\)。
\subsubsection{霍普夫整环}
设$R$是一个整环，$K$是其分式域。若$H$是$K$上的一个霍普夫代数，则$R$上的一个霍普夫整环$O$是$H$ 的一个整环，且在代数与余代数运算下封闭：特别地，余乘法$\Delta : O \to O \otimes O$成立\(^\text{[11]}\)。
\subsubsection{类群元素}
一个类群元素是满足$\Delta(x) = x \otimes x$的非零元素 $x$。类群元素在反元给出的逆元运算下构成一个群\(^\text{[12]}\)。

一个原始元素$x$ 则满足：$\Delta(x) = x \otimes 1 + 1 \otimes x$\(^\text{[13][14]}\)。 
\subsection{例子}
\begin{table}[ht]
\centering
\caption{请输入表格标题}\label{tab_HPFds1}
\begin{tabular}{|c|c|c|c|c|c|c|c|}
\hline
 & \textbf{取决于} & \textbf{余乘法} & \textbf{余元} & \textbf{反元} & \textbf{交换性} & \textbf{余交换性} & \textbf{备注} \\
\hline
群代数 $KG$ & 群 $G$ & $\Delta(g) = g \otimes g \quad \forall g \in G$ & $\varepsilon(g) = 1 \quad \forall g \in G$ & $S(g) = g^{-1} \quad \forall g \in G$ & 当且仅当 $G$ 是阿贝尔群时成立 & 是 &  \\
\hline
从有限群到 $K$ 的函数 $f$，$K^G$（带点加与点乘运算） & 有限群 $G$ & $\Delta(f)(x,y) = f(xy)$ & $\varepsilon(f) = f(1_G)$ & $S(f)(x) = f(x^{-1})$ & 是 & 当且仅当 $G$ 是阿贝尔群时成立 &  \\
\hline
紧群上的表象函数 & 紧群 $G$ & $\Delta(f)(x,y) = f(xy)$ & $\varepsilon(f) = f(1_G)$ & $S(f)(x) = f(x^{-1})$ & 是 & 当且仅当 $G$ 是阿贝尔群时成立 & 反之，每一个在 $\mathbb{C}$ 上具有有限 Haar 积分的交换的、对合的、约化的霍普夫代数都以这种方式出现，这给出了Tannaka–Krein 对偶性的一种表述\(^\text{[15]}\)。 \\
\hline
代数群上的正则函数 &  & $\Delta(f)(x,y) = f(xy)$ & $\varepsilon(f) = f(1_G)$ & $S(f)(x) = f(x^{-1})$ & 是 & 当且仅当 $G$ 是阿贝尔群时成立 & 反之，域上的每一个交换霍普夫代数都以这种方式来自一个群概形，从而给出了范畴之间的一种反等价\(^\text{[16]}\)。 \\
\hline
张量代数 $T(V)$ & 向量空间 $V$ & $\Delta(x) = x \otimes 1 + 1 \otimes x, \; x \in V,\; \Delta(1) = 1 \otimes 1$ & $\varepsilon(x) = 0$ & $S(x) = -x \quad \forall x \in T^1(V)$（并推广到更高阶张量幂） & 当且仅当 $\dim(V) = 0,1$ 时成立 & 是 & 对称代数与外代数（它们是张量代数的商）在此余乘法、余元和反元定义下也都是霍普夫代数 \\
\hline
普遍包络代数 $U(\mathfrak{g})$ & 李代数 $\mathfrak{g}$ & $\Delta(x) = x \otimes 1 + 1 \otimes x \quad \forall x \in \mathfrak{g}$（此规则与对易子相容，因此可以唯一地扩展到整个 $U$） & $\varepsilon(x) = 0 \quad \forall x \in \mathfrak{g}$（同样可扩展到 $U$） & $S(x) = -x$ & 当且仅当 $\mathfrak{g}$ 是阿贝尔李代数时成立 & 是 &  \\
\hline
Sweedler 的霍普夫代数$H = K[c, x] / (c^2 = 1, \; x^2 = 0, \; xc = -cx)$。 & $K$ 是一个特征不等于 2 的域。 & $\Delta(c) = c \otimes c,\;\; \Delta(x) = c \otimes x + x \otimes 1,\;\; \Delta(1) = 1 \otimes 1$ & $\varepsilon(c) = 1,\;\; \varepsilon(x) = 0$ & $S(c) = c^{-1} = c,\;\; S(x) = -cx$ & 否 & 否 & 其底层向量空间由 $\{1, c, x, cx\}$ 张成，因此维数为 4。这是最小的同时非交换且非余交换的霍普夫代数例子。 \\
\hline
对称函数环\(^\text{[17]}\) &  & 以完全齐次对称函数 $h_k \; (k \geq 1)$ 表示$\Delta(h_k) = 1 \otimes h_k + h_1 \otimes h_{k-1} + \cdots + h_{k-1} \otimes h_1 + h_k \otimes 1$ & $\varepsilon(h_k) = 0$ & $S(h_k) = (-1)^k e_k$ & 是 & 是 &  \\
\hline
\end{tabular}
\end{table}
注意：有限群上的函数可以与群环对应起来，尽管它们更自然地被视为对偶结构——群环由有限个元素的和组成，因此可以通过将函数作用在这些求和元素上，与群上的函数配对。
\subsubsection{李群的上同调}
李群$G$的上同调代数（在域$K$上）是一个霍普夫代数：乘法由杯积给出；余乘法则由群乘法$G \times G \to G$所诱导：
$$
H^{*}(G,K) \;\longrightarrow\; H^{*}(G \times G,K) \;\cong\; H^{*}(G,K) \otimes H^{*}(G,K) .~
$$
这一观察实际上正是霍普夫代数概念的来源。利用这种结构，霍普夫证明了关于李群上同调代数的一个结构定理。

\textbf{定理（Hopf）}\(^\text{[18]}\)设$A$是一个有限维、分次交换、分次余交换的霍普夫代数，定义在特征为 0 的域上。那么，作为代数，$A$是一个自由外代数，其生成元均为奇次数。
\subsection{表示论}
设$A$是一个霍普夫代数，$M, N$是$A$-模。则$M \otimes N$也是一个$A$-模，定义为：
$$
a (m \otimes n) := \Delta(a)(m \otimes n) = (a_{1} \otimes a_{2})(m \otimes n) = (a_{1}m \otimes a_{2}n),~
$$
其中$m \in M, \; n \in N,\; \Delta(a) = (a_1, a_2)$。此外，可以将平凡表示定义为基域 $K$，其作用为：
$$
a(m) := \epsilon(a) m, \quad m \in K.~
$$
最后，还可以定义$A$的对偶表示：若$M$是一个$A$-模，且 $M^{*}$ 是其对偶空间，则有：
$$
(af)(m) := f(S(a)m), ~
$$
其中 $f \in M^{*}, \; m \in M$。

$\Delta, \epsilon, S$ 三者之间的关系保证了某些自然的向量空间同态实际上是 $A$-模同态。例如：向量空间的自然同构$M \to M \otimes K$与$M \to K \otimes M$ 同时也是 $A$-模同构；向量空间同态$M^{*} \otimes M \to K,\; f \otimes m \mapsto f(m)$也是 $A$-模同态；但映射$M \otimes M^{*} \to K$并不一定是$A$-模同态。
\subsection{相关概念}
分次霍普夫代数常用于代数拓扑：它们是$H$-空间所有同调或上同调群的直和上自然出现的代数结构。

局部紧量子群是霍普夫代数的推广，并携带拓扑结构。李群上所有连续函数所成的代数就是一个局部紧量子群。

拟霍普夫代数是霍普夫代数的推广，其中余结合律只在“扭曲”意义下成立。它们被用于研究Knizhnik–Zamolodchikov 方程\(^\text{[21]}\)。

乘子霍普夫代数由 Alfons Van Daele 于 1994 年提出\(^\text{[22]}\)，是霍普夫代数的推广：其余乘法是从一个代数（可含幺元，也可无幺元）映射到该代数与自身张量积的乘子代数。

霍普夫群-(余)代数由 V. G. Turaev 于 2000 年提出，也是霍普夫代数的一种推广。
\subsubsection{弱霍普夫代数}
弱霍普夫代数，又称为量子群胚，是霍普夫代数的推广。与霍普夫代数一样，弱霍普夫代数构成一个自对偶的代数类；即：若$H$是一个（弱）霍普夫代数，则其对偶空间$H^*$（即$H$上的线性形式空间）也是一个（弱）霍普夫代数，其中的代数–余代数结构由$H$与其余代数–代数结构的自然配对给出。
\begin{itemize}
\item 通常，弱霍普夫代数$H$被假设为一个有限维的代数和余代数，带有余积$\Delta: H \to H \otimes H$和余元$\epsilon: H \to k$，它满足霍普夫代数的所有公理，除了可能存在 $\Delta(1) \neq 1 \otimes 1$ 或 $\epsilon(ab) \neq \epsilon(a)\epsilon(b)$（某些 $a,b \in H$）。取而代之，要求满足以下条件：
$$
(\Delta(1)\otimes 1)(1\otimes \Delta(1)) \;=\; (1\otimes \Delta(1))(\Delta(1)\otimes 1) \;=\; (\Delta \otimes \mathrm{Id})\Delta(1),~
$$
$$
\epsilon(abc) \;=\; \sum \epsilon(ab_{(1)}) \, \epsilon(b_{(2)}c) \;=\; \sum \epsilon(ab_{(2)}) \, \epsilon(b_{(1)}c),~
$$
对所有$a,b,c \in H$成立。
\item 此外，$H$拥有一个弱化的反元$S: H \to H$，它满足以下公理：\\
1.$S(a_{(1)})a_{(2)} = 1_{(1)} \, \epsilon(a1_{(2)}), \quad \forall a \in H$（右侧是一个投影，通常记作 $\Pi_R(a)$ 或 $\epsilon_s(a)$，其像是一个称为 $H_R$ 或 $H_s$ 的可分子代数）；\\
2.$a_{(1)} S(a_{(2)}) = \epsilon(1_{(1)}a) \, 1_{(2)}, \quad \forall a \in H$（这是另一个投影，通常记作 $\Pi_L(a)$ 或 $\epsilon_t(a)$，其像是一个可分代数 $H_L$ 或 $H_t$，并且通过 $S$ 与 $H_L$ 反同构）；\\
3.$S(a_{(1)})a_{(2)}S(a_{(3)}) = S(a), \quad \forall a \in H$.需要注意的是，如果 $\Delta(1) = 1 \otimes 1$，那么这些条件便化简为霍普夫代数反元的通常两条条件。
\end{itemize}
这些公理部分上是为了保证$H$-模范畴是一个刚性单oidal范畴。其中的单位 $H$-模就是上文提到的可分代数 $H_L$。

例如，有限群胚代数就是一个弱霍普夫代数。特别地，集合$[n]$上的群胚代数，若在其中每对元素$i, j \in [n]$之间都有一对可逆箭头$e_{ij}$与$e_{ji}$，则它同构于$n \times n$矩阵代数$H$。在这个特殊的$H$上，其弱霍普夫代数结构由下式给出：余积：$\Delta(e_{ij}) = e_{ij} \otimes e_{ij}$，余元：$\epsilon(e_{ij}) = 1$，反元：$S(e_{ij}) = e_{ji}$。
在此情形下，可分子代数 $H_L$ 与 $H_R$ 重合，并且它们是非中心的交换代数（即对角矩阵的子代数）。

关于弱霍普夫代数的早期理论性贡献，可以参见\(^\text{[23]}\)和\(^\text{[24]}\)。
\subsubsection{霍普夫代数胚}
\begin{itemize}
\item 参见 Hopf algebroid。
\end{itemize}
\subsubsection{与群的类比}
群可以通过与霍普夫代数相同的图表（或等价地，相同的运算）来公理化，只是此时 $G$被视为一个集合而非模。在这种情况下：
\begin{itemize}
\item 域$K$被替换为一个单点集；
\item 存在一个自然的余元（映射到单点）；
\item 存在一个自然的余乘法（对角映射）；
\item 单位元是群的单位元；
\item 乘法就是群的乘法；
\item 反元就是群的逆元。
\end{itemize}
从这个角度来看，可以把群看作是定义在“一元域”上的一个霍普夫代数\(^\text{[25]}\)。
\subsection{编织单oidal范畴中的霍普夫代数}
霍普夫代数的定义可以自然推广到任意编织单oidal范畴\(^\text{[26][27]}\)。在这样一个范畴$(C, \otimes, I, \alpha, \lambda, \rho, \gamma)$中，一个霍普夫代数是一个六元组$(H, \nabla, \eta, \Delta, \varepsilon, S)$，其中 $H$是$C$中的一个对象，并且：\\
$\nabla : H \otimes H \to H$ —— 乘法；\\
$\eta : I \to H$ —— 单位；\\
$\Delta : H \to H \otimes H$ —— 余乘法；\\
$\varepsilon : H \to I$ —— 余元；\\
$S : H \to H$ —— 反元；

这些都是范畴$C$中的态射，并满足相应的条件。

1）三元组$(H, \nabla, \eta)$是单oidal范畴$(C, \otimes, I,\alpha, \lambda, \rho, \gamma)$中的一个幺半群，即下列图表交换\(^\text{[b]}\)：
\begin{figure}[ht]
\centering
\includegraphics[width=14.25cm]{./figures/738aea6d6792a468.png}
\caption{} \label{fig_HPFds_2}
\end{figure}
2）三元组$(H, \Delta, \varepsilon)$是单oidal范畴$(C, \otimes, I, \alpha, \lambda, \rho, \gamma)$中的一个余幺半群，即下列图表交换\(^\text{[b]}\)：
\begin{figure}[ht]
\centering
\includegraphics[width=14.25cm]{./figures/97f7c0702b2d47cd.png}
\caption{} \label{fig_HPFds_3}
\end{figure}
3）在$H$上的幺半群与余幺半群结构是相容的：乘法$\nabla$与单位$\eta$是余幺半群的态射，并且（在这种情形下是等价的）余乘法$\Delta$与余元 $\varepsilon$ 同时也是幺半群的态射；这意味着下列图表必须交换：
\begin{figure}[ht]
\centering
\includegraphics[width=14.25cm]{./figures/82afc9f712255c6b.png}
\caption{} \label{fig_HPFds_4}
\end{figure}
其中，$\lambda_{I}: I \otimes I \to I$是范畴$C$中的左单位态射，而 $\theta$ 是函子之间的一个自然变换：$(A \otimes B) \otimes (C \otimes D) \;\;\xrightarrow{\;\theta\;}\;\; (A \otimes C) \otimes (B \otimes D)$,它在所有由范畴 $C$ 的结构变换（结合律、左右单位、对换及其逆变换）所构成的自然变换类中是唯一的。

满足 1）、2）、3）性质的五元组$(H, \nabla, \eta, \Delta,\varepsilon)$称为范畴$(C, \otimes, I, \alpha, \lambda, \rho, \gamma)$中的一个双代数。

4）反元的图表是交换的：
\begin{figure}[ht]
\centering
\includegraphics[width=10cm]{./figures/83786c1bcc57ee22.png}
\caption{} \label{fig_HPFds_5}
\end{figure}
典型例子如下：
\begin{itemize}
\item \textbf{群}.在集合范畴$(\text{Set}, \times, 1)$中（其中张量积是笛卡尔积 $\times$，单位对象取为一个任意的单点集，例如 $1=\{\varnothing\}$），三元组$(H, \nabla, \eta)$是范畴意义下的幺半群，当且仅当它是通常代数意义下的幺半群，即运算$\nabla(x,y) = x \cdot y, \quad \eta(1)$的行为与通常的乘法和单位元一致（但元素 $x \in H$ 可能没有逆元）。与此同时，三元组$(H, \Delta, \varepsilon)$是范畴意义下的余幺半群，当且仅当 $\Delta$ 是对角映射：$\Delta(x) = (x,x)$.并且余元 $\varepsilon$ 也能唯一确定：$\varepsilon(x) = \varnothing$。任何这样的余幺半群结构 $(H,\Delta,\varepsilon)$ 都与任意的幺半群结构 $(H,\nabla,\eta)$ 相容，即定义第 3 部分中的图表始终交换。推论：在 $(\text{Set},\times,1)$中，每一个幺半群$(H,\nabla,\eta)$自然可以被看作是一个双代数$(H,\nabla,\eta,\Delta,\varepsilon)$，反之亦然。在这种双代数中，若存在反元$S:H \to H$，则恰好意味着每个元素$x \in H$在乘法$\nabla(x,y)=x \cdot y$ 下都有逆元$x^{-1} \in H$。因此，在集合范畴 $(\text{Set},\times,1)$ 中，霍普夫代数恰好就是通常代数意义下的群。
\item \textbf{经典霍普夫代数}.在特殊情形下，若$(C,\otimes,s,I)$是某个域$K$上的向量空间范畴，那么其中的霍普夫代数正是上述的经典霍普夫代数。
\item \textbf{群上的函数代数}.群上的标准函数代数$\mathcal{C}(G),\; \mathcal{E}(G),\; \mathcal{O}(G),\; \mathcal{P}(G)$（分别表示连续函数、光滑函数、全纯函数和正则函数）是在stereotype 空间范畴$(\text{Ste},\odot)$ 中的霍普夫代数\(^\text{[28]}\)。
\item \textbf{群代数}.群上的 stereotype 群代数$\mathcal{C}^{\star}(G),\; \mathcal{E}^{\star}(G),\; \mathcal{O}^{\star}(G),\; \mathcal{P}^{\star}(G)$（分别表示测度、分布、解析泛函与流）是在 stereotype 空间范畴$(\text{Ste},\circledast)$中的霍普夫代数\(^\text{[28]}\)。这些霍普夫代数被用于非交换群的对偶理论\(^\text{[29]}\)。
\end{itemize}
\subsection{参见}
\begin{itemize}
\item 拟三角霍普夫代数
\item 代数/集合类比
\item 霍普夫代数的表示论
\item 缎带霍普夫代数
\item 超代数
\item 超群
\item 任意子李代数
\item Sweedler 的霍普夫代数
\item 置换的霍普夫代数
\item Milnor–Moore 定理
\end{itemize}
\subsection{注释与参考文献}
\subsubsection{注释}\\
a.群$G$的有限性意味着$KG \otimes KG$ 自然同构于$K[G \times G]$。这被用于上文中的余乘法公式。对于无限群$G$，$KG \otimes KG$是$K[G \times G]$的真子集。在这种情况下，可以赋予有限支撑函数空间一个霍普夫代数结构。\\
b.这里$\alpha_{H,H,H} : (H \otimes H) \otimes H \;\to\; H \otimes (H \otimes H)$,$\lambda_H : I \otimes H \;\to\; H$,$
\rho_H : H \otimes I \;\to\; H$分别是单oidal范畴$(C, \otimes, I, \alpha, \lambda, \rho, \gamma)$中的结合律、左单位和右单位的自然变换。
\subsubsection{引文}
\begin{enumerate}
\item Haldane, F. D. M.; Ha, Z. N. C.; Talstra, J. C.; Bernard, D.; Pasquier, V. (1992). *Yangian symmetry of integrable quantum chains with long-range interactions and a new description of states in conformal field theory. Physical Review Letters. 69 (14): 2021–2025. Bibcode:1992PhRvL..69.2021H. doi:10.1103/physrevlett.69.2021. PMID 10046379.
\item Plefka, J.; Spill, F.; Torrielli, A. (2006). Hopf algebra structure of the AdS/CFT S-matrix. Physical Review D. 74 (6) 066008. arXiv\:hep-th/0608038. Bibcode:2006PhRvD..74f6008P. doi:10.1103/PhysRevD.74.066008. S2CID 2370323.
\item Abreu, Samuel; Britto, Ruth; Duhr, Claude; Gardi, Einan (2017-12-01). Diagrammatic Hopf algebra of cut Feynman integrals: the one-loop case. Journal of High Energy Physics. 2017 (12): 90. arXiv:1704.07931. Bibcode:2017JHEP...12..090A. doi:10.1007/jhep12(2017)090. ISSN 1029-8479. S2CID 54981897.
\item Underwood 2011, p. 55
\item Underwood 2011, p. 62
\item Dăscălescu, Năstăsescu & Raianu (2001). "Prop. 4.2.6". Hopf Algebra: An Introduction. p. 153.
\item Dăscălescu, Năstăsescu & Raianu (2001). "Remarks 4.2.3". Hopf Algebra: An Introduction. p. 151.
\item Quantum groups lecture notes
\item Nichols, Warren D.; Zoeller, M. Bettina (1989). *A Hopf algebra freeness theorem. American Journal of Mathematics, 111 (2): 381–385. doi:10.2307/2374514. JSTOR 2374514. MR 0987762.
\item Montgomery 1993, p. 36
\item Underwood 2011, p. 82
\item Hazewinkel, Michiel; Gubareni, Nadezhda Mikhaĭlovna; Kirichenko, Vladimir V. (2010). *Algebras, Rings, and Modules: Lie Algebras and Hopf Algebras*. Mathematical surveys and monographs. Vol. 168. American Mathematical Society. p. 149. ISBN 978-0-8218-7549-0.
\item Mikhalev, Aleksandr Vasilʹevich; Pilz, Günter, eds. (2002). The Concise Handbook of Algebra. Springer-Verlag. p. 307, C.42. ISBN 978-0-7923-7072-7.
\item Abe, Eiichi (2004). Hopf Algebras. Cambridge Tracts in Mathematics. Vol. 74. Cambridge University Press. p. 59. ISBN 978-0-521-60489-5.
\item Hochschild, G (1965). Structure of Lie groups. Holden-Day, pp. 14–32.
\item Jantzen, Jens Carsten (2003). Representations of algebraic groups. Mathematical Surveys and Monographs, vol. 107 (2nd ed.), Providence, R.I.: American Mathematical Society. ISBN 978-0-8218-3527-2. Section 2.3.
\item Hazewinkel, Michiel (January 2003). Symmetric Functions, Noncommutative Symmetric Functions, and Quasisymmetric Functions. Acta Applicandae Mathematicae. 75 (1–3): 55–83. arXiv\:math/0410468. doi:10.1023/A:1022323609001. S2CID 189899056.
\item Hopf, Heinz (1941). Über die Topologie der Gruppen–Mannigfaltigkeiten und ihre Verallgemeinerungen. Ann. of Math. 2 (in German). 42 (1): 22–52. doi:10.2307/1968985. JSTOR 1968985.
\item Underwood 2011, p. 57
\item Underwood 2011, p. 36
\item Montgomery 1993, p. 203
\item Van Daele, Alfons (1994). *Multiplier Hopf algebras* (PDF). Transactions of the American Mathematical Society. 342 (2): 917–932. doi:10.1090/S0002-9947-1994-1220906-5.
\item Böhm, Gabriella; Nill, Florian; Szlachanyi, Kornel (1999). Weak Hopf Algebras. J. Algebra. 221 (2): 385–438. arXiv\:math/9805116. doi:10.1006/jabr.1999.7984. S2CID 14889155.
\item Nikshych, Dmitri; Vainerman, Leonid (2002). Finite groupoids and their applications. In Montgomery, S.; Schneider, H.-J. (eds.). New directions in Hopf algebras. Vol. 43. Cambridge: M.S.R.I. Publications. pp. 211–262. ISBN 978-0-521-81512-3.
\item Group = Hopf algebra, « Secret Blogging Seminar, Group objects and Hopf algebras », video of Simon Willerton.
\item Turaev & Virelizier 2017, 6.2.
\item Akbarov 2009, p. 482.
\item Akbarov 2003, 10.3.
\item Akbarov 2009.
\end{enumerate}
\subsubsection{参考文献}
\begin{itemize}
\item Dăscălescu, Sorin; Năstăsescu, Constantin; Raianu, Șerban (2001), Hopf Algebras. An Introduction, Pure and Applied Mathematics, vol. 235 (1st ed.), Marcel Dekker, ISBN 978-0-8247-0481-0, Zbl 0962.16026.
\item Cartier, Pierre (2007), "A Primer of Hopf Algebras", in Cartier, P.; Moussa, P.; Julia, B.; Vanhove, P. (eds.), Frontiers in Number Theory, Physics, and Geometry, vol. II, Berlin: Springer, pp. 537–615, doi:10.1007/978-3-540-30308-4\_12, ISBN 978-3-540-30307-7.
\item Fuchs, Jürgen (1992), Affine Lie algebras and quantum groups. An introduction with applications in conformal field theory, Cambridge Monographs on Mathematical Physics, Cambridge: Cambridge University Press, ISBN 978-0-521-48412-1, Zbl 0925.17031.
\item Heinz Hopf, Über die Topologie der Gruppen-Mannigfaltigkeiten und ihrer Verallgemeinerungen, Annals of Mathematics 42 (1941), 22–52. Reprinted in Selecta Heinz Hopf, pp. 119–151, Springer, Berlin (1964). MR 0004784, Zbl 0025.09303.
\item Montgomery, Susan (1993), Hopf algebras and their actions on rings, Regional Conference Series in Mathematics, vol. 82, Providence, Rhode Island: American Mathematical Society, ISBN 978-0-8218-0738-5, Zbl 0793.16029.
\item Street, Ross (2007), Quantum groups: A Path To Current Algebra, Australian Mathematical Society Lecture Series, vol. 19, Cambridge University Press, ISBN 978-0-521-69524-4, MR 2294803, Zbl 1117.16031.
\item Sweedler, Moss E. (1969), Hopf algebras, Mathematics Lecture Note Series, W. A. Benjamin, Inc., New York, ISBN 978-0-8053-9254-8, MR 0252485, Zbl 0194.32901.
\item Underwood, Robert G. (2011), An Introduction to Hopf Algebras, Berlin: Springer-Verlag, ISBN 978-0-387-72765-3, Zbl 1234.16022.
\item Turaev, Vladimir; Virelizier, Alexis (2017), Monoidal Categories and Topological Field Theory, Progress in Mathematics, vol. 322, Springer, doi:10.1007/978-3-319-49834-8, ISBN 978-3-319-49833-1.
\item Akbarov, S.S. (2003). "Pontryagin duality in the theory of topological vector spaces and in topological algebra". Journal of Mathematical Sciences. 113 (2): 179–349. doi:10.1023/A:1020929201133. S2CID 115297067.
\item Akbarov, S.S. (2009). "Holomorphic functions of exponential type and duality for Stein groups with algebraic connected component of identity". Journal of Mathematical Sciences. 162 (4): 459–586. arXiv:0806.3205. doi:10.1007/s10958-009-9646-1. S2CID 115153766.
\end{itemize}
